\paragraph{Summary of the hypotheses.}
\begin{center}\footnotesize
\begin{tabular}{@{}lll@{}}
\toprule
Hypothesis & Result & Evidence \\
\midrule
H1a Jev's belief not worse than card counting & not supported & log loss +0.17 to +0.28 (st.~1), +0.29 to +0.34 (st.~2) \\
H1b card counting in state: too confident & supported & top probability 0.69 vs 0.49 given \\
H1c belief improves later in the game & not supported & Figure~\ref{fig:belief} \\
H2a paper's ranking in A & partly & holds in 240-s games, not in 10,000-event games \\
H2b B: fewer trades, smaller fundamentalist edge & supported & 35 vs 114 trades; 112 vs 222 chips \\
H2c Jev $-$ twin differs between A and B & supported & A $-$ B $= +25$ [12, 36] \\
H3a simple rules obeyed & partly & ``never take'' broken 1--14\% of the time \\
H3b calculations not done & supported & fundamentalist overlap 0.36--0.43 \\
H3c words act literally & supported & behaviour noise trader buys 57\% vs 8\% \\
H3d wording effect same in A and B & not supported & algorithm better in 5/6 types in A, 2/6 in B \\
H4a neutral Jev between f and n & partly & below f; below n in A, not clear in B \\
H4b similar decisions, similar profit & not supported & $\rho(\text{overlap}, |d|) = +0.53$ \\
\bottomrule
\end{tabular}
\end{center}

\paragraph{What the results mean.} A persona text changes what Jev does. It does not make Jev trade like the trader
that the text describes, and it does not make Jev earn what that trader earns. Three results explain this.
\begin{enumerate}
\item \textbf{Jev does not know the goal suit.} In both stages, Jev's belief is worse than a uniform guess, and it does
not improve during the game. The fundamentalist, the only rule-based trader that earns money, depends on this belief.
So the fundamentalist persona is the one that fails most (86 to 106 chips per game behind its twin).
\item \textbf{Jev follows the words, not the procedure.} Jev obeys short rules well and reads descriptions literally.
It does not carry out a calculation that the text describes. A persona is therefore a useful way to steer Jev's
\emph{style} (how often it acts, which side it takes), but not its \emph{judgement}.
\item \textbf{Jev spreads its probability.} We sample from Jev's probabilities over up to 253 options. Jev puts some
probability on many prices, while each twin chooses one price by rule. Sampling therefore makes Jev trade often and at
prices far from the true value. We think that this is the main reason why Jev trades twice as often as its twin and why market prices
become less accurate when Jev plays (we did not test this directly).
\end{enumerate}

\paragraph{Market rules.} The rules change the outcome as much as the persona does. The regression effect of rule set A
($+20$ chips) is larger than any wording effect. In B, every trade clears all quotes, and only an order that beats the
best quote is accepted. The rule-based twins lose many orders to this rule (17 rejected orders per game); Jev does not,
because its menu has only accepted prices. But Jev still earns less in B. In B, the prices of trades with Jev are further from the true value
(9.5 vs 6.2 with the twin; 7.2 vs 5.9 in A), and Jev's belief is worse (log loss 1.64 vs 1.56). We think that these
two facts cause the larger loss in B, but our design does not test this. The rule sets are a design choice with a large effect, and a study of AI
traders should report more than one.

\paragraph{For simulation of market participants.} Our goal was to simulate real market participants with Jev. The
results say: with a text persona only, Jev gives a trader with its own recognisable style, not a copy of a given model
trader. It is closest to the passive types (market makers, noise trader) and furthest from the informed type
(fundamentalist). If a simulation needs informed traders, the information must be part of the design; it cannot come
from the persona text. The +assist result (stage~1) shows that giving Jev the calculated belief is not enough either,
because Jev then becomes too confident.

\paragraph{Limitations.}\label{sec:limits}
\begin{itemize}
\item \textbf{Power.} The noise check gives a minimum detectable difference of 28--43 chips per condition with 40
games. Smaller differences between Jev and its twin, and small wording effects, can stay hidden. (The run stopped once
after 38 games, when the API key reached its spending limit, and continued with the same seeds; all conditions have the
same 40 games.)
\item \textbf{Bottom-feeder persona text.} The bottom-feeder persona says that it copies ``each other player'', but
the twin copies only the fixed fundamentalist opponent (as in the paper). Jev's view of the text is therefore wider than
the twin's rule.
\item \textbf{Homogeneous noise-trader games do not trade.} The paper does not say what a noise trader does in a suit
with no bid. In our reading, it does not trade there, so four noise traders never trade, while the paper reports trades.
This does not change stage~2, where each game has a fundamentalist and a bottom-feeder.
\item \textbf{Card counting is a weak reference.} The paper's calculation uses only the cards seen in trades and the
player's own hand. Its log loss is only slightly below a uniform guess. A stronger reference (for example, a model of the
other players' behaviour) could show an even larger gap to Jev.
\item \textbf{One model, one decoding.} We test one version of Jev (\texttt{typesafe/jev-1.13}) and sample from its full
distribution, as fixed in advance. Choosing Jev's most probable option instead would give different, less random
behaviour (in an earlier pilot, argmax decoding made Jev pass almost always).
\item \textbf{Values read from figures.} The paper's values in Figure~\ref{fig:rep} are read by eye from its figures.
\end{itemize}
